\section{Before Starting the Campaign}

Here are some tips and suggestions intended for the GM. Following them is not mandatory, but it might help you prepare and run the game.

\subsection{Required Knowledge}

As the story starts in Kholinar and spends most of the time there, the characters are expected to be familiar with Alethi culture, and as an extension - the players should also know the basics, especially about alethi fauna, flora, spren and raging highstorms.

In the later chapters, the story will move alongside the events of \bookWOR{} and \bookOB{}. The players will experience the Everstorm, rebirth of the Singers, Revelry of \npcHeartRevel{}, and the dangers of Shadesmar. Make sure you are familiar with those.

\subsection{Read the Adventure}

You, as a GM, are expected to know the overall story of the campaign, so you could present the events with greater ease to your players, and, if a situation calls for it, modify them to your table's needs. Please familiarize yourself with most important NPCs, their goals and beliefs.

\subsection{Players’ Expectations}
\label{ssec:playersExpectations}

Inform the players that the story will start in Kholinar and may contain spoilers from the books, especially \bookWOR{} and \bookOB{}. Later, some events of \bookWAT{} may be also presented. 

Please warn the players that Radiant Spren are rare in the area. Discuss with them beforehand if they are expecting to join any specific Radiant Order and what are the reasons for that. If a player explicitly wishes to play non-Enlightened Radiant, they should be allowed to do so, but the adventure assumes all Radiants are by default the Enlightened ones (\textbf{note:} that doesn't mean the spren and their knight are evil). 

Please make sure you know what spren to introduce for a player aiming for a Radiant Path. The book allows for an easy bonding with Enlightened spren in \nameref{chap:chapter2}, but the characters may meet their spren before that moment (especially if they are aiming for non-Enlightened path).

On top of that, please discuss with your players following topics:
\begin{itemize}
	\item How serious should the adventure be?
	\item What topics players would prefer to avoid during the game?
	\item Is there anything the players would like to see in the game?
	\item Are the players comfortable with their characters dying?
	\item Are the players comfortable with characters’ romance, nudity or other sexual content?
	\item If you are of age - should your group allow alcoholic drinks?
	\item What rule variants should be in effect?
	\item How often should your group meet and how long should a single game session take?
	\item How should you handle situations when one player cannot attend a game session?
\end{itemize}

Gathering answers to those questions is not a requirement, but it will help you create a comfortable and enjoyable atmosphere for everyone.

\textbf{NOTE:} In \nameref{chap:chapter5}, the characters will enter an area influenced by \npcHeartRevel{}, the \npcHeartRevelTitle{}, encountering many NPCs indulging their deepest desires. This book will provide relatively mild and family-friendly descriptions, but if you are aiming for more daring scenes, please make sure all players at the table are comfortable with that.

\textbf{NOTE:} In the end of \nameref{chap:chapter7}, the party will visit a secret hospital in \locationKharbranth{}. Some of the scenes there may be too graphic and gruesome for the more sensitive players. Please make sure you know the boundaries your table is comfortable with, before you proceed with those descriptions.

\subsection{Create Player Characters}

Unless your players have their characters ready after previous adventures, they must create a new party. This book contains several pre-made character sheets you can use in your game (see \hyperref[chap:exampleSheets]{Appendix \ref{chap:exampleSheets}}.), but you should encourage players to create more personalized heroes for the story.

While creating new characters, you might want to follow those guidelines:
\begin{itemize}
	\item Encourage Human Ancestry, unless the player is familiar with more complex Singer rules and culture.
	\item Encourage Alethi culture or at least an expertise in it, as most of the story will take place in that area.
	\item Share \hyperref[backstoryIdeas]{\textbf{Character Backstory Ideas table}} with your players for possible inspirations for characters, their motivations and goals for the story.
	\item If a player plans playing as Knight Radiant, make sure you know what spren to introduce in the story.
	\item Warn players that the usual Radiant Spren may not be easily encountered in the area.
\end{itemize}

\subsection{Character Relationships}

Allow players do determine how their characters know each other and what is the reason they are traveling together. Establishing characters’ relations before the start of the adventure will make it less likely the heroes will be motivated to split the party and follow their own goals.

You can add more details to character’s relationship by suggesting each player to ask about some specific details two other characters. Players can think of their own questions they want to ask each other or they can choose something from the list of below:

\begin{itemize}
	\item I noticed you play with some trinket. What is it and why is it so important to you?
	\item It seems you have some skills not considered appropriate for you by alethi culture. How did you learn them?
	\item We share some unusual hobby/interest. What is it and how did we discover that fact about us?
	\item You shared a secret with me, once. What was it and why did you trust me with it?
	\item When we first met you didn't do a good impression on me. Why, and what happened it changed later?
	\item You once told me that doing the right thing is not always right. What did you mean by that?
	\item Despite living in Kholinar you don't always act like a typical alethi. What are the differences in your lifestyle?
	\item You once showed me you carry a drawing of a spren. What spren is it and why do you have this picture?
	\item You seem to have rather peculiar view on parshmen. What is it and how did you develop it?
	\item We once discussed our views on vorin church. What was your opinion on that matter?
	\item Do you remember the reign of \npcDeadKing{}? How do you find the new king?
	\item What was the most terrible battle that changed you? Why, and how did it change you as a person?
\end{itemize}

\newpage
\subsection{Character Backstory}

Encourage players to write a short character backstory to better root them in the world and the story. Consider if any NPCs from those backstories could make an appearance in this adventure and how they can enrich the story. You can modify scenes described in this book to include those NPCs, for example:
\begin{itemize}
	\item An NPC from the PC's backstory may be met in Kholinar. They might need protection, or escalate the situation, putting others in danger.
	
	\item The party can meet an NPC they know in the hospital of \npcTaravangian{}.
	
	\item Someone they know might have went to war on the Shattered Plains. The party might meet them in Urithiru.
\end{itemize}

You can also invite players to share some connections with the NPCs appearing in the story (\npcTroublemaker{} and \npcGrandma{}). That will make the first chapters of the story much more smoother and personal.

It is not a secret for the citizens, that there are several fractions pulling the strings in the city: \npcQueen{} and her Palace Guard, the Ardents and the Wall Guard. Some even whisper about some mysterious criminal organization operating in the shadows. The PCs would probably know about them and may even have some ties with one of them.